\subsection{Experimental Procedure: Training, Selection and Deployment}
\label{Subsection:ExperimentalProcedure}

An agent design does not by itself produce a result. Training is stochastic, so the same configuration run twice gives two different policies, and a procedure is needed that says which of them is reported and on what grounds. This subsection describes that procedure, the criteria applied at each stage, and the two deployment decisions that sit between the learned policy and the orders actually sent.

\subsubsection{Training Procedure: Seeds, Episodes and Data Splits}

The eleven-year window is divided once into three consecutive periods. Writing the window as $[T_0, T_3)$ with $T_0$ the first of January 2015 and $T_3$ the first of January 2026:

\begin{equation}
\underbrace{[T_0, T_1)}_{\text{training, 9 years}} \quad
\underbrace{[T_1, T_2)}_{\text{validation, 1 year}} \quad
\underbrace{[T_2, T_3)}_{\text{testing, 1 year}}
\label{Equation:Split}
\end{equation}

Weights are updated on the training period only. The validation period decides which state of the network is kept. The testing period is visited once, after every other decision has been taken. Within that split, a single training run proceeds as follows. A seed fixes the network initialisation, the exploration noise stream and the order in which the replay buffer is sampled, and thereby determines the entire run. The agent then performs $30$ episodes, an episode being one pass over the training period. Weights are not reset between episodes: the run is a single trajectory of $30$ passes, and the network at the start of episode $n+1$ is the network at the end of episode $n$. After each episode the agent is evaluated once over the validation period with learning disabled, and that evaluation is reduced to a single number by the Calmar ratio, its annualised return divided by its maximum drawdown. Writing $f_n$ for that value and $g_n$ for an eligibility test defined below, the state retained after episode $n$ is

\begin{equation}
\theta^{\star} \leftarrow \theta_n \quad \text{if} \quad (g_n \wedge \neg g^{\star}) \;\vee\; \left(g_n = g^{\star} \wedge f_n > f^{\star}\right)
\label{Equation:KeepBest}
\end{equation}

which keeps an eligible episode over an ineligible one regardless of score, and otherwise keeps the higher score within the same eligibility class. The saved state is a copy: the run continues training from $\theta_n$, not from $\theta^{\star}$, so a later episode is never restarted from an earlier checkpoint. The eligibility test $g_n$ rejects policies that are degenerate rather than merely poor. An episode qualifies only if it opened at least ten positions, spent at least three hundred bars on each side of the market, and held its weaker side for at least $30\%$ of the time it held its stronger one. A policy that never trades, or that is long for the entire period, can produce a respectable Calmar ratio on a trending window while having learned nothing about when to hold a position, and the test removes it from consideration before the score is compared. Between $128$ and $256$ seeds were run for each currency pair, differing in nothing but the seed. Figure~\ref{Figure:TrainingPipeline} shows how the three stages fit together.

\begin{figure}[htb!]
\centering
\begin{tikzpicture}[
  node distance=5mm,
  stage/.style={draw, rounded corners=2pt, align=center, font=\small, text width=32mm, minimum height=8mm, inner sep=3pt},
  small/.style={stage, font=\scriptsize, text width=30mm},
  note/.style={font=\scriptsize\itshape, align=left},
  flow/.style={-stealth, thick}
]
\node[stage] (seed) {Seed $s$\\\scriptsize init, noise, sampling};
\node[small, right=10mm of seed] (episode) {Episode $n$ of 30\\one pass over training};
\node[small, right=10mm of episode] (valid) {Validation pass\\learning disabled};
\node[small, below=9mm of valid] (keep) {Calmar $f_n$ + eligibility\\keep-best checkpoint};
\node[small, below=9mm of episode] (test) {Test pass\\once, after 30 episodes};
\node[stage, below=9mm of seed] (pool) {Pool of 128--256\\trained seeds};
\node[stage, below=9mm of pool] (gates) {Gates\\reject degenerate};
\node[small, right=10mm of gates] (sweep) {Risk sweep\\$\times 1.0$ to $\times 3.0$};
\node[small, right=10mm of sweep] (score) {Composite score\\percentile ranks};
\node[stage, below=9mm of gates] (champion) {Published model};

\draw[flow] (seed) -- (episode);
\draw[flow] (episode) -- (valid);
\draw[flow] (valid) -- (keep);
\draw[flow] (keep) -- node[note, above, midway] {} (test);
\draw[flow] (keep.west) to[out=180, in=0] node[note, above, midway] {\scriptsize next episode} (episode.south east);
\draw[flow] (test) -- (pool);
\draw[flow] (pool) -- (gates);
\draw[flow] (gates) -- (sweep);
\draw[flow] (sweep) -- (score);
\draw[flow] (score.south) to[out=270, in=0] (champion.east);
\end{tikzpicture}
\caption{The three stages of the procedure. The Calmar ratio chooses among episodes within one seed; the composite score of Section~\ref{Subsection:ExperimentalProcedure} chooses among seeds.}
\label{Figure:TrainingPipeline}
\end{figure}

One property of this arrangement must be stated explicitly, because it limits what the results in Section~\ref{Section:ExperimentalResults} can claim. The split of Equation~\ref{Equation:Split} governs the training of an individual seed. The choice of which seed to publish, described next, is made on performance measured over the whole eleven-year window, and the performance figures reported in this work are measured over that same window. The testing period is therefore held out from the training of each candidate and from the choice among its episodes, but not from the choice among candidates, and the headline returns are not out-of-sample estimates. This is the weakest point in the method and it is not disguised anywhere in what follows.

\subsubsection{Model Selection: Gates and Composite Ranking}

Once the seeds for a pair have been trained, each is evaluated over the full window under the cost model of Table~\ref{Tables:CostModel}, and one is chosen. The choice is made in two steps, rejection followed by ranking, because the two questions are different: whether a policy is admissible at all, and which of the admissible ones is best.

Four gates reject a candidate outright. Its return must be positive. Its weaker direction must account for at least $10\%$ of its exposure, which removes policies that are long or short throughout. Its mean holding period must be at least twenty-four bars, which removes policies whose signal is noise. And its regime score, the fraction of market movement during which it was positioned correctly, must exceed the score its own policy achieves when its position sequence is rotated in time. That last gate deserves a note, because a fixed threshold was tried first and does not work. A rotated policy holds the same positions for the same durations but at the wrong moments, so its score measures what that policy's autocorrelation and holding structure would earn by luck alone. Measured across the seven pairs those nulls range from $48$ to $70$, so a fixed bar of, say, $55$ would be below chance on one pair and needlessly severe on another. The null is therefore computed per model rather than assumed.

Candidates surviving the gates are ranked by a composite of percentile ranks within their own pool. For a metric $m$ and a pool of $K$ candidates, the rank of candidate $i$ is

\begin{equation}
R_m(i) = 100 \cdot \frac{\left|\left\{ j : v_m(j) < v_m(i) \right\}\right|}{K - 1}
\label{Equation:PercentileRank}
\end{equation}

where $v_m$ is negated for metrics that are better when smaller, such as drawdown. The metrics are divided into three groups, listed in Table~\ref{Tables:Composite}, and the score is the unweighted mean of the three group means:

\begin{equation}
S(i) = \frac{1}{3} \sum_{G \in \{\text{money}, \text{safety}, \text{behaviour}\}} \frac{1}{|G|} \sum_{m \in G} R_m(i)
\label{Equation:Composite}
\end{equation}

\begin{table}[htb!]
\centering
\footnotesize
\begin{tabularx}{\textwidth}{@{}lX@{}}
\toprule
\textbf{Group} & \textbf{Metrics, each entering as a percentile rank within the candidate pool} \\
\midrule
Money & Annualised return, Calmar ratio, Sterling ratio \\
\addlinespace
Safety & Sharpe ratio, Sortino ratio, inverse maximum drawdown, inverse downside volatility \\
\addlinespace
Behaviour & Two-sidedness (the weaker exposure side as a fraction of half), regime edge over the model's own rotation null, mean holding period, longest holding period \\
\bottomrule
\end{tabularx}
\caption{Components of the composite selection score of Equation~\ref{Equation:Composite}. The three groups carry equal weight, so return contributes one ninth of the total rather than dominating it. The ratios are those defined in Section~\ref{Subsection:PerformanceMeasurement}.}
\label{Tables:Composite}
\end{table}

Two design decisions are embedded in Equation~\ref{Equation:Composite}. Ranks are relative to the pool rather than absolute, so no threshold can exclude a candidate that is simply better than everything around it. And the three groups carry equal weight, so a policy cannot win on return alone: money, safety and behaviour each contribute a third, and a candidate with the highest return in its pool but no two-sidedness and no regime edge scores poorly on two groups out of three. Each surviving candidate is also evaluated at several position sizes, from the reference risk of Equation~\ref{Equation:PositionSizing} up to three times it. Scaling the risk fraction does not change the learned policy, only the size of the positions it takes, and the level chosen for each pair is reported with the results.

\subsubsection{Deployment: Decision Cadence and Rebalance Threshold}

Two decisions separate the learned policy from the orders that are actually sent, and between them they account for more of the delivered performance than any hyperparameter.

The first is the decision cadence. The agent observes every hourly bar, but it is permitted to act only once per calendar day. A key is formed from the timestamp of each bar and a new action is committed only when that key changes:

\begin{equation}
\kappa(t) = (\text{year}(t), \text{month}(t), \text{day}(t))
\label{Equation:DecisionKey}
\end{equation}

The key is built from the calendar rather than from a count of bars, which makes it robust to the gaps and short sessions described in Section~\ref{Subsection:TradingFramework}: a week containing a holiday still produces one decision per trading day rather than drifting out of alignment. The second is a rebalance threshold. A continuous action changes slightly at every decision, and acting on every such change would send a stream of tiny orders each paying the spread and a commission. An order is therefore sent only when the change in the target position is large enough to be worth its cost:

\begin{equation}
\left| V_{target} - V_{current} \right| \;\geq\; \max\!\left(V_{min},\; \eta \left| V_{ref} \right|\right), \qquad \eta = 0.20
\label{Equation:Hysteresis}
\end{equation}

where $V_{min}$ is the broker's minimum volume. The threshold is a fraction of the reference volume rather than of the position currently held, so it does not vanish as a position is reduced. Table~\ref{Tables:Deployment} isolates the contribution of each decision, measured on one model with nothing retrained. The comparison is between four configurations of the same weights, so the differences are attributable to the deployment layer alone.

\begin{table}[htb!]
\centering
\footnotesize
\begin{tabular}{@{}llrrrr@{}}
\toprule
\textbf{Cadence} & \textbf{Threshold} & \textbf{Commission} & \textbf{Return} & \textbf{Sharpe} & \textbf{Max drawdown} \\
\midrule
Daily & --- & none & $+42.48\%$ & $0.319$ & $24.4\%$ \\
Hourly & --- & none & $-38.57\%$ & $-0.308$ & $50.5\%$ \\
Hourly & --- & 3.5 points & $-64.70\%$ & $-0.732$ & $70.5\%$ \\
Daily & $0.20$ & 3.5 points & $+34.02\%$ & $0.275$ & $23.5\%$ \\
\bottomrule
\end{tabular}
\caption{Contribution of the two deployment decisions, measured on one EUR/USD model over the full window with nothing retrained. All four rows are the same weights under different execution settings, so the differences are attributable to the deployment layer alone. This model precedes the seven-pair campaign and serves here only to isolate the effect.}
\label{Tables:Deployment}
\end{table}

Two things follow from that table. Moving from hourly to daily decisions is worth some eighty percentage points, turning a losing configuration into a profitable one, because an hourly cadence acts on noise the eleven-year horizon of the discount factor was never meant to capture. And the rebalance threshold recovers roughly twelve percentage points under realistic costs while reducing drawdown rather than increasing it, which is unusual: commission alone costs over twenty percentage points without the filter and about six with it. A continuous action space is what makes the strategy expressive, and this filter is what makes it affordable.

\subsubsection{Worked Example: One Position From Decision to Settlement}

The preceding subsections define each component separately. This one follows a single position through all of them, using a trade taken by the delivered EUR/USD model, so that the arithmetic can be checked rather than assumed.

At 02:00 on 30 January 2019 the daily decision key of Equation~\ref{Equation:DecisionKey} changed, and the agent was consulted for the first time that day. It received the 38-value observation of Section~\ref{Subsection:AgentDesign}, built from the bar that had just closed, and returned a positive action. The reference volume of Equation~\ref{Equation:PositionSizing} at that moment, risking $1\%$ of the balance over a stop of $1.5$ average true ranges, gave a target of $12{,}000$ units of the base currency once the action had scaled it and the result was rounded to the broker's volume step. The change from the position then held exceeded the threshold of Equation~\ref{Equation:Hysteresis}, so an order was sent. The tick recorded at that instant quoted $1.14392$ ask against $1.14387$ bid, a spread of five points. The buy was filled at the ask. Twenty-four hours later the key changed again, the agent's action reversed the sign, and the position was closed at the bid then quoted, $1.14931$. The profit follows from those two prices and the volume, and is expressed in the account currency at the conversion rate recorded on the closing tick:

\begin{equation}
12{,}000 \times \left(1.14931 - 1.14392\right) = 64.68 \text{ USD} \;\;\longrightarrow\;\; 56.27 \text{ EUR}
\label{Equation:WorkedGross}
\end{equation}

The spread is already inside that figure. It was paid at entry by filling at the ask rather than the mid, and again at exit by filling at the bid, which is what makes the round trip cross it exactly once. The commission is charged separately at both ends:

\begin{equation}
2 \times 3.5 \times 10^{-5} \times 12{,}000 = 0.84 \text{ USD} \;\;\longrightarrow\;\; 0.72 \text{ EUR}
\label{Equation:WorkedCommission}
\end{equation}

The swap is zero, the position having been held overnight on a swap-free account. The position therefore settled at $56.27 - 0.72 = 55.55$ EUR, which is the figure recorded in the exported trade record. One trade says little. Across all $1{,}777$ trades the delivered EUR/USD model took over the eleven years, the same accounting gives a gross profit of $8{,}733.62$ EUR, a commission charge of $1{,}669.50$ EUR and a swap charge of exactly zero, for a net of $7{,}064.12$ EUR. Commission alone consumed $19.1\%$ of the gross profit, and that is before the spread, which is not separable from the gross figure because it is embedded in every fill. A result reported without these charges would be describing a different strategy.